%! Graphing Quadratic Functions & Transformations — Unit 2, Lesson 2.1 — Guided Notes (Answer Key)
\documentclass[10pt]{article}
\usepackage{algebra2-article}
\usepackage{algebra2-key}   % pulls in -boxes; do NOT also load -boxes
\newcommand{\TallMath}[1]{$\displaystyle #1\rule[-1.4em]{0pt}{3.2em}$}
\newcommand{\lgrid}[4]{%
  \draw[step=1, linegray!40, very thin] (#1,#3) grid (#2,#4);
  \draw[->, semithick, charcoal] (#1,0)--(#2,0) node[below right, font=\scriptsize]{$x$};
  \draw[->, semithick, charcoal] (0,#3)--(0,#4) node[left, font=\scriptsize]{$y$};}
% Vocab answer for the key: bold forest term, then the filled definition on a write-line.
\newcommand{\vocabans}[2]{%
  \par\noindent\textbf{\textcolor{forest}{#1:}}\\[1pt]\ansline{#2}\par}

\begin{document}

\pageheader{Unit 2, Lesson 2.1}{Guided Notes --- Answer Key}

\vspace{0.10in}

\begin{objectivebox}
\textbf{By the end of this lesson, I will be able to\dots}
\begin{itemize}\setlength{\itemsep}{2pt}
  \item read a parabola as \emph{transformations} of the parent $y=x^2$: $f(x)+k$, $f(x-h)$, and $a\,f(x)$;
  \item find the \emph{vertex}, \emph{axis}, \emph{intercepts}, and \emph{max/min} from a quadratic in \emph{vertex}, \emph{standard}, or \emph{intercept} form;
  \item use those features to \emph{graph} a quadratic or match it to its graph, and convert between forms.
\end{itemize}
\end{objectivebox}

\vspace{0.06in}

\begin{vocabbox}
\small
Fill in each term as we build it together.
\par\vspace{2pt}
\vocabans{Vertex form}{$f(x)=a(x-h)^2+k$; the vertex is $(h,k)$ and the axis of symmetry is $x=h$}
\vocabans{Standard form}{$f(x)=ax^2+bx+c$; the $y$-intercept is $(0,c)$ and the axis is $x=-\frac{b}{2a}$}
\vocabans{Intercept (factored) form}{$f(x)=a(x-p)(x-q)$; the zeros are $x=p$ and $x=q$, and the axis is halfway between them}
\vocabans{Horizontal shift $f(x-h)$}{slides the graph right if $h>0$, left if $h<0$ --- opposite the sign inside the parentheses}
\vocabans{Vertical stretch / reflection $a\,f(x)$}{narrower if $|a|>1$, wider if $0<|a|<1$; reflects over the $x$-axis if $a<0$}
\end{vocabbox}

\begin{hookbox}
These three equations are all \textbf{the same parabola}:
\[ f(x)=x^2-2x-3 \qquad f(x)=(x-1)^2-4 \qquad f(x)=(x+1)(x-3). \]
\begin{itemize}\setlength{\itemsep}{1pt}
  \item Which one shows the \textbf{vertex} at a glance? \ansline{$(x-1)^2-4$ --- vertex form, vertex $(1,-4)$.}
  \item Which one shows the \textbf{zeros} at a glance? \ansline{$(x+1)(x-3)$ --- intercept form, zeros $x=-1$ and $x=3$.}
\end{itemize}
Each form is cut to show off a different feature. Today we read all three --- and use them to graph.
\end{hookbox}

\boxguard[22]
\begin{notesbox}{1. The Parent $y=x^2$ and Its Transformations}
Every parabola is the parent $y=x^2$ after some \textbf{shifts}, \textbf{stretches}, and
\textbf{reflections}. Track the parent's anchor point $(2,4)$ and its vertex $(0,0)$ through each move.
\begin{center}
\begin{tikzpicture}[scale=0.42]
  % Parent
  \begin{scope}
    \lgrid{-3}{3}{-1}{5}
    \begin{scope}\clip (-3,-1) rectangle (3,5);
      \draw[very thick, forest, domain=-2.2:2.2, samples=60, smooth] plot (\x,{(\x)*(\x)});
    \end{scope}
    \fill[charcoal] (0,0) circle (3pt) node[above left, font=\scriptsize]{$(0,0)$};
    \node[charcoal, font=\scriptsize] at (0,-2){parent $y=x^2$};
  \end{scope}
  % Shifted g(x) = (x-1)^2 - 4
  \begin{scope}[xshift=10cm]
    \lgrid{-3}{5}{-5}{3}
    \draw[navy, dashed, semithick] (1,-5)--(1,3);
    \begin{scope}\clip (-3,-5) rectangle (5,3);
      \draw[very thick, forest, domain=-1.2:3.2, samples=70, smooth] plot (\x,{(\x-1)*(\x-1) - 4});
    \end{scope}
    \fill[charcoal] (1,-4) circle (3pt) node[below right, font=\scriptsize]{$(1,-4)$};
    \node[charcoal, font=\scriptsize] at (1,-6.4){$g(x)=(x-1)^2-4$};
  \end{scope}
\end{tikzpicture}
\end{center}
Fill in the effect of each change and where the vertex $(0,0)$ lands:
\begin{center}
\renewcommand{\arraystretch}{1.5}
\begin{tabularx}{\linewidth}{@{}l X l@{}}
\textbf{Change to $y=x^2$} & \textbf{Effect on the graph} & \textbf{Vertex $\to$} \\
\hline
$f(x)+k$ & shifts \ans{up} ($k>0$) or down ($k<0$) & $(0,\ \textcolor{keyred}{\mathbf{k}})$ \\
$f(x-h)$ & shifts \ans{right} $h$ units (careful: sign flips!) & $(\textcolor{keyred}{\mathbf{h}},\ 0)$ \\
$a\,f(x)$, $a>0$ & \ans{narrower} if $a>1$, wider if $0<a<1$ & $(0,0)$ \\
$-f(x)$ & reflects over the \ans{$x$}-axis (opens down) & $(0,0)$ \\
\end{tabularx}
\end{center}
Put them together: the graph of $\;a(x-h)^2+k\;$ is the parent moved so its vertex is at
$(\textcolor{keyred}{\mathbf{h}},\ \textcolor{keyred}{\mathbf{k}})$.
\end{notesbox}

\begin{notesbox}{2. Vertex Form: $f(x)=a(x-h)^2+k$}
Vertex form \emph{is} the transformations written down. Read the vertex right off it. For
$f(x) = (x-1)^2 - 4$:
\begin{center}
\begin{tikzpicture}[scale=0.5]
  \lgrid{-3}{5}{-5}{3}
  \draw[navy, dashed, semithick] (1,-5)--(1,3);
  \node[navy, font=\scriptsize, right] at (1.05,2.4){axis $x=1$};
  \begin{scope}\clip (-3,-5) rectangle (5,3);
    \draw[very thick, forest, domain=-1.4:3.4, samples=75, smooth] plot (\x,{(\x-1)*(\x-1) - 4});
  \end{scope}
  \fill[charcoal] (1,-4) circle (3pt) node[below right, font=\scriptsize]{vertex $(1,-4)$};
  \fill[charcoal] (0,-3) circle (3pt) node[left, font=\scriptsize]{$(0,-3)$};
  \fill[charcoal] (-1,0) circle (3pt) node[above left, font=\scriptsize]{$(-1,0)$};
  \fill[charcoal] (3,0) circle (3pt) node[above right, font=\scriptsize]{$(3,0)$};
\end{tikzpicture}
\end{center}
\begin{enumerate}[leftmargin=1.6em, itemsep=3pt]
  \item Match to $a(x-h)^2+k$: $a=\textcolor{keyred}{\mathbf{1}}$, $h=\textcolor{keyred}{\mathbf{1}}$, $k=\textcolor{keyred}{\mathbf{-4}}$.
  \item \textbf{Vertex} $(h,k) = (\textcolor{keyred}{\mathbf{1}},\ \textcolor{keyred}{\mathbf{-4}})$; \textbf{axis of symmetry} $x=\textcolor{keyred}{\mathbf{1}}$.
  \item Since $a=1>0$, it opens \ans{up}, so the vertex is a \textbf{minimum} with min value $\textcolor{keyred}{\mathbf{-4}}$.
  \item \textbf{Check it is the same curve.} Expand: $(x-1)^2-4 = \ans{$x^2-2x-3$}$ --- standard form!
\end{enumerate}
\end{notesbox}

\boxguard[22]
\begin{notesbox}{3. Standard Form: $f(x)=ax^2+bx+c$}
Here the vertex is hidden. Two facts unlock it: the $y$-intercept is $(0,c)$, and the axis of symmetry
is $x=-\dfrac{b}{2a}$. For $f(x)=x^2-6x+5$ ($a=1,\ b=-6,\ c=5$):
\begin{center}
\begin{minipage}{0.52\linewidth}
\begin{enumerate}[leftmargin=1.6em, itemsep=4pt]
  \item \textbf{Axis:} $x=-\dfrac{b}{2a}=-\dfrac{-6}{2(1)}=\textcolor{keyred}{\mathbf{3}}$.
  \item \textbf{Vertex:} substitute that $x$: $f(3)=(3)^2-6(3)+5=\textcolor{keyred}{\mathbf{-4}}$, so the vertex is $(\textcolor{keyred}{\mathbf{3}},\ \textcolor{keyred}{\mathbf{-4}})$.
  \item \textbf{$y$-intercept:} $(0,\ \textcolor{keyred}{\mathbf{5}})$.
  \item Opens \ans{up} ($a>0$); range \ans{$y\ge -4$}.
\end{enumerate}
\end{minipage}%
\begin{minipage}{0.46\linewidth}\centering
\begin{tikzpicture}[scale=0.45]
  \lgrid{-1}{7}{-5}{6}
  \draw[navy, dashed, semithick] (3,-5)--(3,6);
  \begin{scope}\clip (-1,-5) rectangle (7,6);
    \draw[very thick, forest, domain=-0.2:6.2, samples=80, smooth] plot (\x,{(\x)*(\x) - 6*(\x) + 5});
  \end{scope}
  \fill[charcoal] (3,-4) circle (3pt);
  \fill[charcoal] (0,5) circle (3pt);
  \fill[charcoal] (1,0) circle (3pt);
  \fill[charcoal] (5,0) circle (3pt);
\end{tikzpicture}
\end{minipage}
\end{center}
\emph{Why $-\tfrac{b}{2a}$?} The axis sits \emph{halfway} between any two symmetric points --- including
the two zeros --- so it is their average. That midpoint is exactly $-\tfrac{b}{2a}$.
\end{notesbox}

\begin{practicebox}
\textbf{A. Match each vertex-form equation to its graph} by reading the vertex $(h,k)$ and the opening
direction. Write A, B, or C.
\begin{center}
\begin{tikzpicture}[scale=0.38]
  % A: (x-2)^2 - 1  vertex (2,-1) up
  \begin{scope}
    \lgrid{-1}{5}{-2}{4}
    \begin{scope}\clip (-1,-2) rectangle (5,4);
      \draw[very thick, forest, domain=0.3:3.7, samples=60, smooth] plot (\x,{(\x-2)*(\x-2) - 1});
    \end{scope}
    \fill[charcoal] (2,-1) circle (2.5pt);
    \node[charcoal, font=\scriptsize] at (2,-3.2){\textbf{A}};
  \end{scope}
  % B: -(x-1)^2 + 3  vertex (1,3) down
  \begin{scope}[xshift=8cm]
    \lgrid{-2}{4}{-2}{4}
    \begin{scope}\clip (-2,-2) rectangle (4,4);
      \draw[very thick, forest, domain=-1.2:3.2, samples=60, smooth] plot (\x,{-1*(\x-1)*(\x-1) + 3});
    \end{scope}
    \fill[charcoal] (1,3) circle (2.5pt);
    \node[charcoal, font=\scriptsize] at (1,-3.2){\textbf{B}};
  \end{scope}
  % C: (x+2)^2 + 1  vertex (-2,1) up
  \begin{scope}[xshift=16cm]
    \lgrid{-5}{1}{-1}{6}
    \begin{scope}\clip (-5,-1) rectangle (1,6);
      \draw[very thick, forest, domain=-4.2:-0.2, samples=60, smooth] plot (\x,{(\x+2)*(\x+2) + 1});
    \end{scope}
    \fill[charcoal] (-2,1) circle (2.5pt);
    \node[charcoal, font=\scriptsize] at (-2,-2.2){\textbf{C}};
  \end{scope}
\end{tikzpicture}
\end{center}
\begin{enumerate}[leftmargin=1.9em, itemsep=3pt]
  \item $y=(x-2)^2-1 \Rightarrow$ vertex \ans{$(2,-1)$}, opens \ans{up}, graph \ans{A}
  \item $y=-(x-1)^2+3 \Rightarrow$ vertex \ans{$(1,3)$}, opens \ans{down}, graph \ans{B}
  \item $y=(x+2)^2+1 \Rightarrow$ vertex \ans{$(-2,1)$}, opens \ans{up}, graph \ans{C}
\end{enumerate}
\textbf{B. Intercept form.} For $q(x)=(x+1)(x-3)$: zeros $x=\ans{$-1$}$ and $x=\ans{$3$}$; axis
$x=\ans{$1$}$; vertex $\ans{$(1,-4)$}$.
\end{practicebox}

\end{document}
